% для очистки бибера - clean up corrupted biber cash: rm -rf $(biber --cache)
% LuaLaTeX, 
% Engine: biber
%%%%%%%%%%%%%%%%%% Document Setup %%%%%%%%
\documentclass[10pt]{article}
\usepackage{filecontents}
\usepackage{xpatch}
\usepackage{graphicx}
\usepackage{metalogo}
\usepackage{amsmath}
\usepackage{textcomp}  % numero symbol
\usepackage{gensymb} % degree symbol
\usepackage[dvipsnames,svgnames,x11names,xcdraw,table]{xcolor}
\usepackage{lmodern} % latin modern
\usepackage[T2A,T1]{fontenc}
\usepackage[utf8]{inputenc}
\usepackage[greek,ngerman,russian,french,english,]{babel}

% ------------------------ для разных буллитов в списках -------------------------------
\usepackage{enumitem} 
\usepackage{pifont}
% ----------------------------------------------------------------------------

\usepackage{booktabs} %%% таблица
\usepackage{tabularx}
\usepackage{array,etoolbox}
\usepackage{boldline}
\usepackage{colortbl} %% цветная таблица
	\setlength{\arrayrulewidth}{0.3mm}
%%%%%% нумерация внутри таблицы %%%

\preto\tabular{\setcounter{magicrownumbers}{0}}
\newcounter{magicrownumbers}
\newcommand\rownumber{\stepcounter{magicrownumbers}\arabic{magicrownumbers}}

\usepackage{longtable}

\usepackage{ifthen}
\usepackage{luacode}
\usepackage{tikz} % load xcolor before tikz
\usetikzlibrary{graphs,graphdrawing}

\usepackage{siunitx}
\usepackage[style=british]{csquotes}
\usepackage{url}
%\makeatletter
\makeatother
\reversemarginpar
\usepackage[paper=letterpaper,
            marginparwidth=0.75in, % Length of section titles
            marginparsep=.05in, % Space between titles and text
            margin=0.8in, % 1 inch margins
            includemp]{geometry}
\setlength{\parindent}{0in}
\usepackage{enumerate} 
\usepackage{etaremune}
\usepackage{bbding}
\usepackage{paralist}

\usepackage{graphics}
\usepackage{wrapfig}
\usepackage{titlesec}
\renewcommand\thesubsection{}
\renewcommand\thesubsubsection{}

\setcounter{secnumdepth}{2}
\usepackage[super]{nth}

%%%%%%%% для цветного параграфа %%%%%%%%
\usepackage{tcolorbox}
%%%%%%%% для цветного параграфа %%%%%%%%

%%%%%%%%%Bib%%%%%%%%%%%%
\usepackage[
  	backend=biber,
	style=numeric,
% 	style=phys,
% 	style=chem-biochem, nature, science(no paper's name): no DOI apa
%	style=, chem-acs, numeric, chem-angew,
  	giveninits=true,
	isbn=true,
	url=true,
	natbib=true,
	defernumbers=true,
	sorting=ndymdt, 
%	sorting=ydnt, % Sort by year (descending), name, title.
	bibencoding=utf8,
	language=auto,    % get main language from babel
	autolang=other,
]{biblatex} 

\usepackage[colorlinks = true,
	breaklinks = true,
	urlcolor = NavyBlue,
	breaklinks = true]{hyperref}

%%%%%%% начало: footnote (\cfoot{abc}) %%
\usepackage{fancyhdr}
\pagestyle{fancy}
\fancyhf{}% to clear existing header/footer if you don't want it
\renewcommand\headrulewidth{0pt}
%%%%%%% конец:  footnote (\cfoot{abc}) %%

%%%%%%% начало: обратная нумерация в библиографии %%%%%%%
\AtDataInput{%
  \ifcsundef{bbx@processedentries:\therefsection}
    {\csgdef{bbx@processedentries:\therefsection}{}}
    {}%
  \xifinlistcs{\thefield{entrykey}}{bbx@processedentries:\therefsection}{}{%
    \listcsxadd{bbx@processedentries:\therefsection}{\thefield{entrykey}}%
    \csnumgdef{bbx@entrycount:\therefsection:\thefield{keywords}}{%
      \csuse{bbx@entrycount:\therefsection:\thefield{keywords}}+1}}}

\DeclareFieldFormat{labelnumber}{\mkbibdesc{#1}}    
\newrobustcmd*{\mkbibdesc}[1]{%
  \number\numexpr\csuse{bbx@entrycount:\therefsection:\thefield{keywords}}+1-#1\relax}
%%%%%%% конец: обратная нумерация в библиографии %%%%%%%
   
\addbibresource{Lemenkova.bib}

%%%%%%%%%%% сортировка библиографии : начало %%%%%%%
\DeclareSortingTemplate{ndymdt}{
  \sort{
    \field{presort}
  }
  \sort[final]{
    \field{sortkey}
  }
  \sort[direction=descending]{
    \field{sortyear}
    \field{year}
    \literal{9999}
  }
  \sort[direction=descending]{
    \field[padside=left,padwidth=2,padchar=0]{month}
    \literal{99}
  }
  \sort[direction=descending]{
    \field[padside=left,padwidth=2,padchar=0]{day}
    \literal{99}
  }
}
%%%%%%%%%%% сортировка библиографии : конец %%%%%%%
 
  %%%%%%%%%%% pagetotal : начало %%%%%%%%%%%
\renewbibmacro*{addendum+pubstate}{%
  \printfield{addendum}%
  \newunit\newblock
  \printfield{pubstate}
  \newunit\newblock
  \ifentrytype{misc}{\printfield{pagetotal}}
  \newunit\newblock
  \ifentrytype{inproceedings}{\printfield{pagetotal}}{}}
 %%%%%%%%%%% pagetotal : конец %%%%%%%%%%%

%%%%%% добавление французского архива HAL L'archive ouverte : начало %%%%%%%
\DeclareFieldFormat{hal}{%
  \mkbibacro{HAL}\addcolon\space
  \ifhyperref
    {\href{https://hal.archives-ouvertes.fr/#1}{\nolinkurl{#1}}}
    {\nolinkurl{#1}}}       
\DeclareFieldAlias{eprint:hal}{hal}
\DeclareFieldAlias{eprint:HAL}{eprint:hal}

\renewbibmacro*{eprint}{%
  \printfield{hal}%
  \newunit\newblock    
  \iffieldundef{eprinttype}
    {\printfield{eprint}}
    {\printfield[eprint:\strfield{eprinttype}]{eprint}}}
%%%%%% добавление французского архива HAL L'archive ouverte : конец %%%%%%%   

\DeclareFieldFormat{url}{%
  \mkbibacro{URL}\addcolon\space
  \href{#1}{\nolinkurl{\thefield{urlraw}}}}

%%%%%%%%%%%% значок Open Access: начало %%%%%%%
\usepackage[tikzsymbol=plos]{biblatex-ext-oa}
\usepackage{kantlipsum}
\DeclareOpenAccessEprintUrl[always]{hal}{%
  http://hal.archives-ouvertes.fr/\thefield{eprint}}
\DeclareOpenAccessEprintAlias{HAL}{hal}
%%%%%%%%%%% значок Open Access:: конец %%%%%%%

%%%%%%%%%%% добавление 2-го doi : начало %%%%%%%%%%%%%
\renewcommand*{\relatedpunct}{\addcolon\space}
\renewcommand*{\relateddelim}{\addcomma\space}

\newbibmacro*{related:prelim}[1]{%
  \renewcommand*{\newunitpunct}{\addcomma\space}%
  \entrydata{#1}{\usebibmacro{doi+eprint+url}}}

\NewBibliographyString{prelim,prelims}
\DefineBibliographyStrings{american}{%
  prelim = {Preliminary version},
  prelims = {Preliminary versions}}
%%%%%%%%%%% добавление 2-го doi : конец %%%%%%%%%%%%%

\usepackage[object=vectorian]{pgfornament}

%%%%%%%%%%%%%%%%%%%%%%%% End Document Setup %%%%%%%

%%%%%%%%%%%%%%%%%%%%%%%%%%% Helper Commands %%%%%%%%%

\newcommand{\makeheading}[1]%
        {\hspace*{-\marginparsep minus \marginparwidth}%
         \begin{minipage}[t]{\textwidth+\marginparwidth+\marginparsep}%
                {\large \bfseries #1}\\[-0.15\baselineskip]%
                 \rule{\columnwidth}{1pt}%
         \end{minipage}}
\renewcommand{\section}[2]%
        {\pagebreak[2]\vspace{1.3\baselineskip}%
         \phantomsection\addcontentsline{toc}{section}{#1}%
         \hspace{0in}%
         \marginpar{
         \raggedright \scshape #1}#2}
        
\newenvironment{innerlist}[1][\enskip\textbullet]%an itemize-style list with little space between items
        {\begin{compactitem}[#1]}{\end{compactitem}}
\newcommand{\blankline}{\quad\pagebreak[2]}
\providecommand*\email[1]{\href{mailto:#1}{#1}}

%%%%%%%%%%%%%%
\usepackage{fontspec}
\usepackage{fontawesome}
\usepackage{wasysym}
\usepackage[misc]{ifsym}
\usepackage{bclogo}
\usepackage{academicons}
\definecolor{orcidlogocol}{HTML}{A6CE39}
\definecolor{orcidlogocol}{HTML}{A6CE39}

%%%%%%%%% End Helper Commands %%%%%%%%%%%%%%%%%
%%%%%%%%%%% Begin CV Document of Polina Lemenkova %%%%%%

\begin{document}

\makeheading {Polina Lemenkova}
\cfoot{\hspace{330pt}{\small{CV of Polina Lemenkova \\ \hspace{320pt}Update: \today}}}
\lfoot{Page \thepage}
%\hspace{1em}

\section{Social Data}
\begin{wrapfigure}{r}{42mm}
%	\begin{center}
	 	\includegraphics[width=0.20\textwidth]{Lemenkova2018.jpg}
%	\end{center}
\end{wrapfigure}

\begin{innerlist}
	\item[\Letter] \email{pauline.lemenkova@gmail.com} \phone +32 471 86 04 59
	%(+007)916-298-3719
	\item[\faHome] Moscow, Russia. \faSkype \space polina.lemenkova
	%\item[\faChild] Moscow, Russia, 1980.
	%\item[\faInfoCircle] Online Profiles: 
	\item[\textcolor{orcidlogocol}{\aiOrcid}] \href{https://orcid.org/0000-0002-5759-1089}{ORCiD ID: 0000-0002-5759-1089},  \faFacebookSquare \space\href{https://www.facebook.com/polina.lemenkova.7}{Facebook}%, \href{https://search.crossref.org/?q=Lemenkova}{Crossref}
	\item[\aiAcademia] \href{https://ulb.academia.edu/PolinaLemenkova}{Academia.edu}, \href{https://ssrn.com/author=3266317}{SSRN}, \href{https://search.datacite.org/people/0000-0002-5759-1089}{DataCite}, \href{https://doi.pangaea.de/10.1594/PANGAEA.680792}{PANGAEA}, \href{https://www.lifescience.net/profile/polina-lemenkova-9skfemem/}{LifeScienceNetwork}
	\item[\faLinkedin] \href{https://www.linkedin.com/in/paulinelemenkova}{Linkedin}, \faSlideshare\space \href{https://www.slideshare.net/PolinaLemenkova/presentations}{SlideShare}, \aiPublons \space \href{https://publons.com/researcher/1587436/polina-lemenkova/}{Publons}, \href{https://elibrary.ru/author_items.asp?authorid=153168}{eLibrary}, \href{https://sciprofiles.com/profile/535013}{SciProfiles}
	\item[\aiFigshare] \href{https://figshare.com/authors/Polina_Lemenkova/5855411}{Figshare}, \href{https://cv.archives-ouvertes.fr/polina-lemenkova}{IdHAL},  \aiarXiv \space\href{https://arxiv.org/a/lemenkova_p_1.html}{arXiv}. \href{http://istina.msu.ru/profile/polinalemenkova/}{IstinaResearcherID. Score: 10,80}. 
	\item[\faGithub] \href{https://github.com/paulinelemenkova}{GitHub}, \href{https://zenodo.org/search?page=1&size=20&q=lemenkova}{Zenodo}, \aiGoogleScholar\space \href{https://scholar.google.com/citations?user=As7apVQAAAAJ&hl=en}{Google Scholar} (Citations: 2161; h-index: 25). 
	\item[\aiResearchGate] \href{https://www.researchgate.net/profile/Polina-Lemenkova-2}{ResearchGate} (Citations: 2,720; RG Score: 48.49).
%	\item[] 
	%(Citations: 1690, \textit{h}-index: 23, \textit{i}10-index: 53). 
%\aiMendeley \space\href{https://www.mendeley.com/profiles/polina-lemenkova/}{Mendeley}, \aiDataverse \href{https://dataverse.harvard.edu/dataverse/harvard?q=lemenkova}{Harvard Dataverse}. %\aiOSF  \href{https://osf.io/zd5w4/}{OSF}, 
%\href{https://www.growkudos.com/profile/287718}{Kudos}, \href{http://dergipark.gov.tr/@lemenkovapolina}{DergiPark},
%\faTwitter \space \href{https://twitter.com/PLemenkova}{Twitter}, \href{https://www.kaggle.com/polinalemenkova}{Kaggle}, \aiPubMed \href{https://www.ncbi.nlm.nih.gov/myncbi/pauline.lemenkova@gmail.com/cv/298055/}{PubMed}, \href{https://www.bibsonomy.org/cv/user/polinalemenkova}{BibSonomy}, \href{http://sciencewise.info/users/Polina_Lemenkova}{ScienceWISE}, \faInstagram \space \href{https://www.instagram.com/polina.lemenkova/}{Instagram} \href{https://profiles.impactstory.org/u/0000-0002-5759-1089}{Impactstory}, \aiOverleaf\space \href{https://www.overleaf.com/read/fxyfmqnzsqqw}{Overleaf \LaTeX}, \aiZotero\space\href{https://www.zotero.org/polina_lemenkova}{Zotero}, \href{https://www.mysciencework.com/home/publications}{MyScienceWork}, \aiSemanticScholar \space \href{https://www.semanticscholar.org/author/Polina-Lemenkova/52194775}{Semantic Scholar}, \href{https://loop.frontiersin.org/people/624535/overview}{Loop} [\aiSciRate] \href{https://scirate.com/polina-lemenkova}{SciRate}, \href{https://www.scopus.com/authid/detail.uri?authorId=55625974200}{Scopus ID: 55625974200} 
% 20+ years of research experience in geoinformatics since 1999 (study, jobs, research stays).
\end{innerlist}

\section{Summary}
\begin{tcolorbox}[text width=14.5cm,colback=FloralWhite,colframe=AntiqueWhite4,]
	\begin{innerlist}
		\item[\Checkmark] \textbf{Education}: M.Sc. in GIS and Earth Observation; B.Sc. in Geography and Cartography.
		\item[\Checkmark] \textbf{Publications}: \emph{80+} journal articles; \emph{90} conference papers; \emph{50} presentations and posters
		\item[\Checkmark] \textbf{Reviewer}: 20+ journal manuscripts (indexed in Scopus/WoS).
		\item[\Checkmark] \textbf{Concepts}: Big Earth Data, Machine Learning, Deep Learning, Data Science, Geomatics
		\item[\Checkmark] \textbf{Skills}: 
			\begin{innerlist}
				\item \textbf{Cartography}: GMT, GRASS GIS, ArcGIS, QGIS, Erdas, ILWIS, Idrisi, ENVI. 
				\item \textbf{Programming/Markup Languages}: R, Python, Octave/MATLAB, \TeX , HTML
				\item \textbf{Unix}: utilities, shell/bash scripts, awk. \href{https://github.com/paulinelemenkova}{GitHub: 100+ repositories}
			\end{innerlist}
		%\item[\Checkmark] \textbf{International profile}: 7 years spent in Europe, 4 years in China; born in Russia
		\item[\Checkmark] \textbf{Languages}: English, Russian. \emph{Also learned}: German, French, Chinese, Spanish, Italian.
	\end{innerlist}
\end{tcolorbox}

\section{Education}

\href{https://www.ulb.be/}{\textbf{Université libre de Bruxelles}} \hfill{\textbf{Belgium, Bruxelles}}\\
09.2021-ongoing: PhD 'Doctorat en Sciences de l'ingénieur et technologie'. \href{https://polytech.ulb.be/en}{École polytechnique de Bruxelles}: Exploitation of historical analog seismological record by image processing and ML. 
%Supervisor: Prof. Dr. Olivier Debeir; Co-supervisor: Dr. Thomas Lecocq.
\vspace{0.7em}

\href{http://www.utwente.nl/en}{\textbf{Universiteit Twente}} \hfill{\textbf{Netherlands, Enschede}}\\
09.2009-03.2011: MSc Course 'GIS and EO for Environmental Modelling and Management'.
\href{http://www.itc.nl/}{Faculty of Geo-Information Science \& Earth Observation}. MSc Thesis: \href{https://zenodo.org/record/2290319}{\emph{'Seagrass mapping and monitoring along the coasts of Crete, Greece'}}. \href{https://drive.google.com/file/d/15nvWlDbIy9DfKZHMaGbFfrdFQcHNEWZS/view?usp=sharing}{MSc Diploma}
\vspace{0.7em}

\href{http://www.uw.edu.pl/en/}{\textbf{University of Warsaw}} \hfill{\textbf{Poland, Warsaw}}\\
07/06-27/07/2010. Environmental Management and Policy in Central Europe. Fieldwork (MSc).
\vspace{0.7em}	

\href{http://www.lunduniversity.lu.se/}{\textbf{Lunds Universitet}}
 \hfill{\textbf{Sweden, Lund}}\\
01-06.2010: GIS Centre. System Dynamics \& Analysis; Advanced GIS and Remote Sensing; Distributed Environmental Modelling; Biodiversity.
\vspace{0.6em}

\href{http://www.southampton.ac.uk/}{\textbf{University of Southampton}} \hfill{\textbf{UK, England}}\\
 09.2009-01.2010. \href{http://www.southampton.ac.uk/geography/}{School of Geography}. \textit {Modules}: GIS for Environmental Modelling; Core Skills in GIS \& RS; Calibration \& Validation of EO Data; Risk Assessment; Physical Geography in Environmental Management. \href{https://drive.google.com/file/d/1I2DASP_waCDPbkQqRW479U-8ERztAKJP/view?usp=sharing}{MSc Diploma}\vspace{0.6em}

\href{http://www.unibas.ch/}{\textbf{Universit\"at Basel}} \hfill{\textbf{Switzerland, Basel}}\\
2008/2009. \href{http://www.physiogeo.unibas.ch/index.php?lang=en}{Geographisches Institut}. \href{https://drive.google.com/file/d/1IT99J8PqoiQZ5xbmlxdN5ty2eNkY6GIi/view?usp=sharing}{ESKAS grantee}. \href{https://www.researchgate.net/project/2008-2009-Switzerland-Basel-ArcGIS-Spatial-Analysis-to-Study-Gallocanta-Lake-Environment}{ArcGIS based spatial analysis}
\vspace{0.6em}

\href{http://www.msu.ru/en/}{\textbf{Lomonosov Moscow State University (MSU)}}
\hfill{\textbf{Russia, Moscow}}\\
01/09/1999 - 25/06/2004. B.Sc. Course at \href{http://www.eng.geogr.msu.ru/}{Faculty of Geography, Dept Cartography \& Geoinformatics}. Graduated with Honors, Summa Cum Laude. \href{https://drive.google.com/file/d/1oRoBCuPVjT7Z39Qq1cgCY2smB2p3qEQI/view?usp=sharing}{Licentiate Diploma: Apostille}. \\
Licentiate Thesis: \emph{\href{https://www.researchgate.net/publication/327955267_Ecological_Mapping_of_Barents_Pechora_Seas}{Ecological Mapping of the Barents and the Pechora Seas}} %\\
%Term Paper: \emph{\href{https://figshare.com/articles/Term_Paper_2002_Ecological_and_Geographical_Mapping_of_the_Baltic_Sea_Region_in_the_Gulf_of_Finland/7294661}{Ecological and Geographical Mapping of the Baltic Sea Region in the Gulf of Finland.}} 

\section{Awards}
	\begin{innerlist}
		\item[\faTrophy]\textbf{University Gold Medal} for Academic Excellence (2004), \href{http://www.msu.ru/en/}{Lomonosov Moscow State University}
	\end{innerlist}

%\pagebreak

\section{Additional Education}

\href{http://eweb.ouc.edu.cn/}{\textbf{Ocean University of China}}, \href{http://eweb.ouc.edu.cn/mgs/}{College of Marine Geo-sciences} \hfill{\textbf{China, Qingdao}}\\
09.2016-07.2020. \textit{ECTS (19)}: Progress in Marine Geology; Methods of Seabed Survey \& Digital Technologies; Continental Margin Tectonics \& Ocean Floor Geodynamics; Sedimentology \& Analysis of Marine Sediments; Advanced Marine Geology; China Studies; Academic Writing. \href{https://drive.google.com/file/d/1DUD1m_ewERWUCYB5G4O3V8-4VD_0AlIg/view?usp=sharing}{Student IDs}.
 \href{https://drive.google.com/file/d/1wWYGR-o37QmaFmTOUsTeBPPdS_kUfhAC/view?usp=sharing}{Transcripts}. %Research Output: \href{http://ssrn.com/abstract=3858289}{Deep-Sea Trenches of the Pacific Ocean: a Comparative Analysis of the Submarine Geomorphology Using Data Modeling by GMT, Python and R}
\vspace{3pt}

\textbf{Charles University in Prague}\hfill{\textbf{Czech Republic, Prague}}\\
2013-2015: Department of Ecology. \textit{3 ECTS}: Environmental Informatics. \href{https://www.researchgate.net/project/2013-2015-Czech-Republic-Sumava-National-Park-Environmental-Analysis}{Šumava National Park}.
\vspace{3pt}

\href{http://www.ut.ee/et}{\textbf{University of Tartu}}\hfill{\textbf{Estonia, Tartu}}\\ 01/05-30/06/2012: \href{https://www.geograafia.ut.ee/et}{Chair of Geoinformatics \& Cartography}. \href{https://www.researchgate.net/project/2012-Estonia-Tartu-IDRISI-GIS-and-CORINE-for-Mapping-Land-Cover-Changes}{Idrisi GIS, CORINE, Landsat TM}.
\vspace{3pt}

\textbf{Technische Universit\"at Dresden} \hfill{\textbf{Germany, Dresden}}\\	
2011-2013. Institute for Cartography. \href{https://www.researchgate.net/project/2011-2013-Germany-TU-Dresden-Environmental-Analysis-of-Tian-Shan-Region-Central-Asia}{Environmental analysis of Tian Shan mountains} 
\vspace{3pt}

\href{http://www.univie.ac.at/}{\textbf{Universit\"at Wien}}, \href{http://www.univie.ac.at/geographie/} {Institut f\"ur Geographie und Regionalforschung} \hfill{\textbf{Austria, Vienna}}\\
01/03-30/06/2011. \href{https://drive.google.com/file/d/1TH5GgB_lV5K1aJVwp3_9aYaqqZMgd7op/view?usp=sharing}{6 ECTS: Risk Assessment \& Disaster Management} \vspace{3pt}

%\textbf{Universit\'{e} de Fribourg}
%\hfill{\textbf{Switzerland, Fribourg}}\\
%01.09.2009-31.03.2010. Facult\'{e} de Sciences. Research focus: Alpine geoarchaeology. \vspace{3pt}

\href{http://www.msu.ru/en/}{\textbf{Lomonosov Moscow State University}}
\hfill{\textbf{Russia, Moscow}}\\
09.2005-06.2007. \href{http://www.msu.ru/en/info/struct/depts/teach.html}{Faculty of Educational Studies}. BSc+ \emph{Using ArcGIS in Teaching Geoscience}.\vspace{3pt}

\href{http://www.msu.ru/en/}{\textbf{Lomonosov Moscow State University}}
\hfill{\textbf{Russia, Moscow}}\\
2004-2007. \href{http://www.eng.geogr.msu.ru/}{Department of Cartography and Geoinformatics}. Modules: Cartography, Philosophy, German. Internship: Teacher of Geodesy at student summer fieldwork, MSU Faculty of Soil Science.

\section{Internships}

\href{http://www.joanneum.at/}{\textbf{Joanneum Research}}. Austria, Graz. 06-09.2003. \href{https://drive.google.com/file/d/1_Vq-Imz1m52B0RyDVdXRN7SLwd0xEDDX/view?usp=sharing}{Student trainee}. Institute of Digital Image Processing. Geo-referencing images, digitizing, calculating, mapping changes in glaciers.

\blankline

\textbf{MPA Kartografija} Russia, Moscow. 06-08.2002. \textit{Student trainee}. Digitizing and editing maps.

%\subsection{Secondary Education}

%\href{http://en.fa.ru}{Financial University under the Government of the Russian Federation}. Moscow. Faculty of Pre-University Education, 1997/98. Modules: Geography, Mathematics, Russian. \href{https://drive.google.com/open?id=1j-EEdq3BZhmYRObJKbr_D4Qay7m5Ypv_}{Certiicate}. \vspace{3pt}
	
%College 'Business Studies in Russia'. Naro-Fominsk. 21.06.1997. Modules: Market Economy, Management System, HR, Enterprise, Management, Marketing, Business Strategy. \href{https://drive.google.com/open?id=1FDlutVuajli_xepJbWY68nK_B9z9a9C-}{Certificate}.\vspace{3pt}

%School \textnumero 4, General Secondary Education. 1987-1997. Russia, Naro-Fominsk. Modules: Russian, Literature, English, Chemistry, Physics, Astronomy, Labor, Physical education, Biology,  Geometry,  Analysis, Informatics, Computers, Geography, Algebra, History, Civilizations. \href{https://drive.google.com/open?id=14yIv-eGwJuK8c2Kafuo7lobs9uObpd6U}{Certificate}.\vspace{3pt}

%School of \faMusic \space \textnumero 47, Piano Faculty. 1987-1995. Russia, Moscow. Modules: Music Literature, Musical Ensemble, Musical Accompaniment, Singing in Choir, Piano, Solfeggio. \href{https://drive.google.com/open?id=1NyfdsM4fF138LzukbnfTDUROBSPuknVr}{Certificate}. 

\section{Technical Skills}

\begin{innerlist}

\item[\faCode\space] \textbf{Programming languages} (codes available on \href{https://github.com/paulinelemenkova}{GitHub}. User of Jupyter Notebook.
	\begin{innerlist}
		\item Python (libraries: Matplotlib, StatsModels, Seaborn, Pandas, NumPy, SciPy, Bokeh). 
		\item R (packages: ggplot2, tidyverse, latticeExtra, dplyr, ggpubr, magrittr, reshape2, etc.)
		\item Octave/MATLAB
		\item Unix shell/bash scripts, Unix utilities, awk, sed, grep
 %(ca 40): at, banner, cat, chmod, cd, cp, date, diff, echo, exit, grep, head, hostname, jobs, kill, less, ls, man, mkdir, mktemp, mv, paste, pg, ps, pwd, rm, sed, seq, sort, sudo, tail, tee, test, touch, w, wc, which, who, whoami.	
%		\item Brief acquaintance with PHP, JavaScript, D3.js.
%		\item Code editors: XCode, Atom. Brief acquaintance with GNU nano, Vim.
%		\item Package managers \& installing systems: MacPorts, pip, npm, Homebrew, Conda (Miniconda3).
	\end{innerlist}
\vspace{1ex}

\item[\faPencil\space] \textbf{Markup languages and typesetting}
\begin{innerlist}
	\item \TeX, HTML/CSS, XML, LyX, gedit, Notepad++
%	\item Editors: \TeX Shop, \TeX Works
	\item Compilers: \LuaLaTeX, \XeLaTeX, pdf\LaTeX, \LaTeX
	\item Bibliographies: Bib\LaTeX, Biber, Bib\TeX, MakeIndex
%	\item \TeX \space distributions: \TeX Live, Mac\TeX, MiK\TeX
\end{innerlist}
	
\item[\faLineChart\space] \textbf{Statistical analysis}: Python, R, Gretl, SPSS, Gnuplot (graph plotting).
	
\item[\faEye\space] \textbf{Cartography and Geoinformatics}: 
\begin{innerlist}
	\item GMT, GRASS GIS, GDAL, QGIS, ArcGIS, Erdas Imagine, ILWIS, ENVI, Idrisi, eCognition
	\item Hydrography: Caris HIPS, ATLAS Hydrosweep DS-2, WASI (Water Color Simulator)
\end{innerlist}

%\item[\faDatabase\space] \textbf{Databases and Data Analysis}
%\begin{innerlist}
%	\item awk (data filtering, sorting, parsing, transforming, aggregation)
%	\item MongoDB (basics)
%	\item OpenRefine (data formatting and filtering)	
%\end{innerlist}
%General skills: data management, administration, organizing into projects, SQL queries.
	
%\item[\faKeyboardO\space] \textbf{OS}: \faApple \space MacOS X. Previously: \faLinux \space (RedHat, Suse), \faWindows \space (XP, Vista, NT, 2000)
%	 \vspace{1ex}
 
%\item[\faPaintBrush\space] \textbf{Graphic Design}:\\
%	PGF/TikZ (\LaTeX), Gimp, Inkscape, CorelDraw, Adobe Photoshop, Adobe Illustrator
	
\end{innerlist}

\section{Languages} 
\begin{innerlist}
	\item[*] English: \href{https://drive.google.com/file/d/14IyQIcHMe95waVXFHXnJpllXY2y1hwQz/view?usp=sharing}{\textsc{IELTS}} C1; British Council, Vienna, Austria, 2011. Score: 7.5 (R8.0, S7.5, W7.5, L6.5).	
	\item[*] German: \href{https://drive.google.com/file/d/1TeuAaleYz2kHNkhR0EzAa4o0Ty4GCZov/view?usp=sharing}{\textsc{Diplom \"OSD}}: C2. Moskauer Staatliche Linguistische Universit\"at, 15.06.2006
	\item[*] French: \href{https://drive.google.com/file/d/13q7YMKsqLFOB-uD0Ao7dYmGosQ_5XaKV/view?usp=sharing}{\textsc{Dipl\^ome DALF}}: C1. Institut Fran\c cais de Br\^eme, Allemagne, 2008.
	\item[*] Chinese: \href{https://drive.google.com/file/d/1lC_0fSOB2GBgjHjqpI1Y1MVf_rfME-Ti/view?usp=sharing}{\textsc{Certificate HSK-5}}. 2.500 characters. Confucius Institute, Qingdao, China, 15.07.2017. 
	\item [*]Italian: \href{https://drive.google.com/file/d/1kv1r05YPYx12nL93V1nSJ8-s9UDlAyWF/view?usp=sharing}{\textsc{Certificato CILS}}: B1, at Instituto Italiano di Cultura di Mosca, 2007.
	\item[*] Spanish: \href{https://drive.google.com/file/d/14j7dgoqT4q8sH-WCjJ32yxLxirUzCEBb/view?usp=sharing}{\textsc{Diploma DELE}} Nivel B1. 22.08.2008, Instituto Cervantes, Mosc\' u, Rusia.
	\item[*] Russian: native.
\end{innerlist}
	
\pagebreak

\section{Grants}

\begin{innerlist}
	\item	(China) China Scholarship Council, Marine Scholarship of China, State Ocean Administration (SOA), 2016SOA002. 09/2016-07/2020. 
	\item (France) Travel grant TRAIL: Formation et recherche pour l'interpr\'etation arch\'eologique des donn\'ees LiDAR. 26-28/03/ 2014 (Besançon-Frasne, Franche-Comté)
	\item (Czechia) Charles University Scholarship
	\item (Netherlands) Groningen Energy Summer School, Rijksuniversiteit Groningen. 06/2013
	\item (China) Taiwan Ministry of Education, Short Term Research Award
	\item (Turkey) T\"{U}B\.{I}TAK. Fellowship 2216. 01/11-31/12/2012.
	\item (Belgium) Bourse d'excellence IN WBI (Wallonie-Bruxelles Int'l). 01/09/-30/10/2012
	\item (Slovakia) N\'{a}rodn\'{y} \v{s}tipendijn\'{y} program SR. SAIA. 01/07- 31/08/2012
	\item (Estonia) DoRa, European Social Fund, Program Activity 5. 01/05-30/06/2012
	\item (Hungary) Hungarian State Scholarship, Balassi Int\'{e}zet (Budapest, HSB). 01/01-30/04/2012
	\item (Finland) CIMO Finnish Centre for International Mobility. 01/07-31/12/2011
	\item (Austria) Stipendiumstiftung \"{O}sterreichische Austauschdienst GmbH (\"{O}AD). 03-06/2011 
	\item Erasmus Mundus Scholarship for joint MSc course. 09/2009-03/2011
	\item (Finland) Finnish Centre for International Mobility CIMO. 01/06-31/08/2009
	\item (Switzerland). ESKAS Swiss Government Excellence Scholarships for Foreign Scholars \& Artists. 15/09/2008-15/06/2009
	\item (Germany) Deutscher Akademischer Austauschdienst DAAD. 10/2007-03/2008
	\item (Russia) Moscow State Lomonosov University Scholarship. 09/2005-06/2007
	\item (Germany) AWI grant for Antarctic expedition ANT/XXIII-4. 02-04.2006
	\item (Austria) Joanneum Research Internship, Inst. of Digital Image Processing. 01/06-31/08/2003
	\item (Russia) Moscow State Lomonosov University Scholarship. 09/1999-06/2004
\end{innerlist}	

%\section{Fieldwork}

%\subsection{Geography Fieldwork}

%\begin{innerlist} 
%	\item Estonia, P\"arnu. May 2012. Observations of the land use types in the H\"a\"ademeeste area
%	\item Greece, Crete, Agia Pelagia. Sep-Oct 2010. Data collection. Seagrass monitoring.
%	\item Poland, Sudety Mts., Krkono\v{s}e, Karpacz (Silesia). July 2010. Environmental monitoring.
%	\item Poland, Bia\l owie\.{z}a National Reserve Forest. July 2010. Monitoring land use types. 
%	\item Poland, Biebrza National Park, Masurian Lakeland (Gi\.zycko, Grajewo). June 2010. Wetlands.
%	\item Antarctica, Amundsen and Bellingshausen Seas. Feb-Apr 2006. Sonar data capture: acoustic deep-sea multibeam echo-sounder Hydrosweep DS; CARIS HIPS; ArcGIS.
%	\item Russia, Chashnikovo. July 2005. Teacher of Geodesy. MSU, Faculty of Soil Sciences. 
%	\item Russia, Gelendzhik. July 2001. Detecting land cover types on aerial images.
%	\item Russia, Khibiny Mts. June 2001. Geodetic surveying (theodolite, tacheometry, GPS).
%	\item Russia, Satino. Jun-Sep 2000. Geomorphology, geodetic measurements: theodolite, GPS; topographic mapping, land use, soil, meteorological \& hydrological measurements.
%\end{innerlist}
	
%\subsection{Archaeological Fieldwork}
%\begin{innerlist} 
%	\item  Cyprus Island, Paphos, Kissonerga. 31.07-14.08.2013. \href{http://www.harparchaeology.co.uk/home}{Heritage \& Archaeological Research Practice}. \href{http://www.harparchaeology.co.uk/field-schools/experimental-archaeology-beer-production-1}{Experimental archaeology: Mediterranean Middle East, Bronze Age}. Survey: dumpy level, drawing elevations, offsetting. Artifacts excavations, sampling, sieving, recording.
	%shovel, brush, pickaxe, bucket, organizing
%	\item Russia, Vybuty Churchyard. 09-13 July 2013. \href{http://arheologpskov.ru/}{Archaeological Center of Pskov Region}. \href{http://www.estlatrus.eu}{Estonia-Latvia-Russia cross border cooperation program}. Artifacts excavation and sorting.
	%site clearing, gridding, shoring, digging, leveling, cleaning artifacts. Settlement AD X-XVII.
%\end{innerlist}

%\pagebreak

\section{Events} 

\begin{etaremune} 

	\item Webinar \emph{Reunión virtual Comisión de Geofísica} on 11th June, 2021, organised by \emph{\href{https://www.ipgh.org/}{Instituto Panamericano de Geografía e Historia}}. \emph{Attendee}.
	\item Webinar \href{https://geospatialworld.net/virtual-events/bentley/smartcities4.0.html}{‘Smart and Sustainable: How Leading Cities Are Achieving Economic, Social and Environmental Sustainability Goals’} on 9th June, 2021 11 AM EDT, organised by Bentley Systems and hosted by Geospatial World. \emph{Attendee}.
	\item	American Geophysical Union (AGU) 24-28 May 2021 (Online). \href{https://agu.confex.com/agu/21workshop2/meetingapp.cgi/Home/0}{Bringing Land, Ocean, Atmosphere and Ionosphere Data to the Community for Hazard Alerts}	
	\item	Russia, Moscow 24.05.2021. \href{https://www.ifz.ru/novosti/24-maya-2021-goda-sostoitsya-seminar-laboratorii-regionalnoj-geofiziki-i-prirodnyix-katastrof}{Seminar of the Laboratory of Regional Geophysics and Natural Disasters}		
	\item	Russia, Moscow 05.11.2020. \href{https://drive.google.com/file/d/14e58LWu4wbZRniu8l9HB7-q6EqJmlDbl/view?usp=sharing}{Web seminar: Information platform Web of Science \& RSCI}		
	\item	Russia, Moscow 31.08.2019. \href{https://corp.mail.ru/ru/press/events/mdsm_aug19/}{Moscow Data Science Major} 		
	\item	Russia, Moscow 29.01-02.02.2019 \nth{51} Tectonics Meeting. Problems of Tectonics of the Continents and Oceans 	
    	\item	Russia, Tambov 14-16.11.2018 Virtual Modeling, Prototyping \& Industrial Design. \href{https://drive.google.com/file/d/1MYXw3WsHmUjq03rlvg3Ti8Gwg1BUu4dr/view?usp=sharing}{Certificate}		
	\item	Russia, Far East, Vladivoskok 19-20.04.2016. Fisheries and Aquaculture. \href{https://drive.google.com/open?id=1-r6f_m2p2iHR98SMeXwXt_M_0WIOYgqU}{Certificate}
	\item Kazakhstan, Almaty 11-12.04.2016 \nth{5} Forum 'Green Bridge Through Generations'. \href{https://drive.google.com/open?id=1s08GHSZrQcXYmJvbcFvF2v68FWi92Lb0}{Certificate}	
	\item	 Russia, Tambov 17-19.11.2015 \nth{2} Conference Virtual Modeling, Prototyping \& Industrial Design. \href{https://drive.google.com/open?id=1d0ijaq5tmPIw48Ma_Sb6EPgbEYF026FZ}{Certificate}
	\item Germany  8-9.10.15 \emph{Workshop DFG Research Training Group Baltic Transcoast, Uni. Rostock} 				
	\item	Russia, Ryazan 29.09.2015 Modern Trends \& Perspectives in Development of Tourism \& Hospitality \href{https://drive.google.com/open?id=1_Fb8oftTRlwbtY5fNpQQCcHTt2QB0-ms}{Certificate} 			
	\item	Russia, Komi, Syktyvkar 14-19.09.2015 \nth{6} Forest Soils Fundumental \& Practical Questions of Forest Soils. \href{https://drive.google.com/open?id=1Hbban7h80uYvQqQmvn3fZI_9v9dhkn_F}{Certificate}. 	
	\item	 Russia, Kazan (Tatarstan) 7-11.07.2015 \nth{13} RSEE Theory \& Practice of the Economic Management of Land Use and Environmental Protection. \href{https://drive.google.com/open?id=1aaiHDa4YTEGarJK400cDoLA7VS92n9pj}{Certificate} 
	\item	Georgia, Kutaisi 3-4.07.2015 \nth{7} Internet \& Society (INSO2015). \href{https://drive.google.com/open?id=1Nu2VPmWyTtNxJsUcSHBoASXjGWw4-DHL}{Invitation Letter} 						
	\item	Russia, Grozny (Chechnya) 19-21.06.2015. Geoenergy. \href{https://drive.google.com/open?id=1cQ-Azw7_42aGD-8XSRH6EgYOwIoQiVI-}{Certificate} 						
	\item	Belarus, Vitebsk 3-4.06.2015. \href{http://conference.vstu.by/}{Conference 'Trends and Perspectives in Regional Systems of Adult Education'}. Plenary Speaker.			
	\item	Belarus, Grodno 28-29.05.2015. Prospects for the Higher School Development \href{https://drive.google.com/open?id=1lEGm8YZTiz3TBunBndk1QtIQft32K0au}{Certificate} 					
	\item	Russia, Ufa (Bashkortostan) 20-22.05.2015. \nth{3} IT-Days. Problems \& Solutions. \href{https://drive.google.com/open?id=1K7bFgsGKMVQ7jyYOg6bA82IyXbsTx_ZO}{Certificate} 					
	\item	Kazakhstan, Aqt\"{o}be 13.05.2015 \href{http://krmu.kz/en/}{\emph{Integration of Education and Science: Challenges of the Modern World}}. \href{https://drive.google.com/open?id=1NxuiGyvuYseBmsPUSfrCwyrgo4WQYG9u}{Certificate}. \href{https://drive.google.com/open?id=11rtqEFLuKOn9M_J_k0Gykv3IY_GBqxFC}{Invitation Letter}. Plenary Speaker.					
	\item	Belarus, Gomel 23-24.04.2015 Geographic Aspects of Sustainable Develop. of Regions. \href{https://drive.google.com/open?id=1KB9rrG8zxeyfXQC04vTdpQYkBoXwwHnI}{Certif.}					
	\item	Russia, Derbent (Dagestan) 15-16.04.2015 \nth{2} \href{http://konferencii.ru/info/111030}{Russian National Conference 'Education Activities in the Mode of Innovation: Concepts, Approaches, Technologies}. Plenary Speaker. 					
	\item	Russia, Belgorod 6-10.04.2015 \href{http://ggf.bsu.edu.ru/Conferences/Conf_04_2015/conf_04_2015.htm}{\nth{3} Int'l Research Conference of Young Scientists \emph{'Geoecology and Sustainable Use of Mineral Resources: Science to Practice}}. Plenary Speaker. \href{https://drive.google.com/open?id=1LQ-unvt7Tb3JGqkEUTGxdzNCumt4Zte_}{Certificate} 				
	\item	Russia, Arzamas 26-27.03.2015. \href{http://matematikum.ucoz.ru/}{Conference on Web-Technologies in the Educational Space: Problems, Approaches, Perspectives}. \href{https://drive.google.com/open?id=1Gmj1xyHSv8_Q2Ej6QTSKjPqhn9nFiIUf}{Certificate}. Plenary Speaker.				
	\item	Russia, Yoshkar-Ola 19.03.2015 \href{http://www.volgatech.net/announcements/konferentsii/79482/}{Actual Problems of State \& Management of Water Resources}				
	\item	Russia, Yekaterinburg (Ural) 27-28.02.2015 \nth{75} \href{http://urgau.ru/index.php/75let-urgau/programma}{Ural State Agrarian University Anniversary}		
	\item Russia, Moscow 21.02.2015. International Open Data Day Moscow.	
	\item	Russia, Voronezh 12-13.02.15 \nth{15} Conf. \href{http://www.cs.vsu.ru/conf/}{Informatics: Problems, Methodologies, Technologies} 
	\item	Russia, Voronezh 12-13.02.2015 \href{http://www.cs.vsu.ru/conf/}{\nth{6} Winter School 'Informatics in Education'} 		
	\item	Czechia, Prague 27.01.2015 \nth{6} \href{https://www.natur.cuni.cz/fakulta/zivotni-prostredi/aktuality/prilohy-a-obrazky/konference/pgs-koference-2015-program}{Annual PGS Conference}		
	\item Hungary, Debrecen 11.12.2014 \emph{\href{http://www.eurisy.org/event-spatial-planning-and-geospatial-services-for-regional-development_26/about}{Workshop 'Spatial Planning \& GIS Services for Regional Development'}}, Interreg. IVC Proj. 'Regions4GreenGrowth', \'{E}szak-Alf\"{o}ld Reg. Development Agency
	\item Greece, Patras 26-28.11.2014. IRLA2014, \nth{1} International Symposium on the Effects of Irrigation \& Drainage on Rural \& Urban Landscapes.			
	\item	Armenia, Tsaghkadzor 17-19.11.2014. \nth{3} \href{http://www.geocom.am/Conference/index_conference.htm}{GIS \& Remote Sensing} 		
	\item Greece, Thessaloniki 22-24.10.14 \emph{\nth{10} International Congress of the Hellenic Geographical Society}. \href{https://drive.google.com/open?id=1ZW460zRhml4pWcQEAz36Lb4AyLkP1O9i}{Certificate}.   				
	\item	Belarus, Minsk 14-17.10.2014 \nth{5} Modern Problems of Geoecology \& Landscapes Studies Belarusian State University, Faculty of Geography \href{https://drive.google.com/open?id=1SV9r3sCw_gkSZAX9OFMTQUXSoxdgrPa-}{Invitation Letter}		
	\item	Greece, Athens 11-13.09.2014. \href{http://canepalconference.wordpress.com/}{CANEPAL International Conference 'Celebrating Pastoral Life; Heritage \& Economic Development}. Plenary Speaker.	
	\item Italy, Rome (Frascati) 20-22.05.2014 \emph{\href{http://seom.esa.int/S2forScience2014/index.php}{SENTINEL-2 for Science ESA Workshop}}			
	\item	Czechia, Olomouc 3-6.04.14. Socio-Economic Transition of China: Opportunities and Threats. \href{https://drive.google.com/open?id=1xCcyg_7t_UR_7LI2OVs-3W5JN1RfTvDb}{Certificate}.							
	\item	Czechia, Brno 06.02.2014. CzechGlobe Conference.
	\item France, Frasne 26-28.03.2014 \emph{TRAIL 2014, Formation et recherche pour l'interpr\'etation arch\'eologique des donn\'ees LiDAR,  Franche-Comt\'{e}}. \href{https://drive.google.com/open?id=1-Tx8qCawFPVtZpjThc3SghWGhBQK-7X-}{Certificate}.
	\item	Croatia, Split 22-24.01.2014 \nth{10} International Conference on Environmental, Cultural, Economic \& Social Sustainability. \href{https://drive.google.com/open?id=1kcPlN3MRb57JhAFF9op1RuKXZRRubVPP}{Certificate}.		
	\item	Spain, Barcelona 3-5.12.2013. Learning Leadership: 'Creating Innovative Learning Communities' (Discussion), 'Challenge of School Design' (Master class). \href{https://drive.google.com/open?id=14TQdhnsq0cSoblaNBmMc3iB7OO7Qp3kY}{Certificate}.	
	\item Portugal, Peniche. 27-29.11.2013. \emph{\nth{6} International Tourism Congress (ITC-13)}. \href{https://drive.google.com/open?id=1py8pVp4P_pybuK_My5rMUa45GJNnCJk8}{Certificate}.			
	\item	Latvia, Riga 14-16.10.2013 Riga Technical University \nth{54} International Scientific Conference. Subsection "Environmental \& Climate Technologies". \href{https://drive.google.com/open?id=1wKzXy2Mkmx9O5Z7-TsEEDEQE7xJ3Rj4T}{Certificate}	
	\item Lithuania, Vilnius 02-05.07.2013 \nth{4} International Meeting of Fire Effects on Soil Properties. \href{https://drive.google.com/open?id=12OOyk3bfJ1HIbNfI3coNDKMzijxvm9i_}{Certificate}.
	\item Netherlands, Enschede 26.06.2013. Participatory GIS Symposium.	
	\item Netherlands 17-28.06.2013 \emph{\href{http://www.rug.nl/research/energy/education/summerschool/}{Groningen Energy Summer School: A multi-disciplinary approach to energy transition, from policy to physics}}. Rijksuniversiteit Groningen. 5 ECTS. \href{https://drive.google.com/open?id=1UqWf3XBSRSZitaddfX8P5P39CCz4lYXS}{Certificate}.	
	\item China, Taipei. 22/05/2013. Seminar at \href{https://www.ntu.edu.tw/}{National Taiwan University (NTU)}. Contributed Talk: 'Using \emph{K-means} algorithm classifier for urban landscapes classification in Taipei area, Taiwan'.
	\item China, Taipei (Taiwan). 10/05/2013. Seminar at \href{https://www.ntu.edu.tw/}{NTU}. Contributed Talk: 'Russia. Moscow. Society, culture, geography'.
	\item Belgium 03/10/2012. Universit\'e Libre de Bruxelles. Contributed Talk: 'Urban sprawl in Estonia'. Joint seminars on Geospatial Analysis, Ecology \& Epidemiology.      				
	\item	Serbia, Belgrade 29-31.05.2012. \nth{3} Int'l Geoscience Student Conference. Chairperson during the Section 'Environmental issues'  \href{https://drive.google.com/open?id=1a5QN7buCbt5-wT231Z0mjrGlvU2TOHa0}{Invitation Letter}		
	\item	Serbia, Belgrade 27-29.05.2012 \nth{3} Int'l Conference Geoscience \& Environment \href{https://drive.google.com/open?id=1a5QN7buCbt5-wT231Z0mjrGlvU2TOHa0}{Invitation}		
	\item	Ukraine, Kiev 14-17.05.2012 \nth{11} \href{http://geoinformatics.org.ua/eng/conferences/}{Geoinformatics - Theoretical \& Applied Aspects}   	
	\item Austria, Vienna 03-08.04.2011 \href{http://meetings.copernicus.org/egu2011/}{European Geoscience Union, General Assembly 2011}

\end{etaremune}	

%\pagebreak

\section{Work History}

%\subsection{Job Employment}

\href{https://polytech.ulb.be/}{\textbf{Université Libre de Bruxelles}}. École polytechnique de Bruxelles (Brussels Faculty of Engineering). Belgium, Bruxelles. 08.2021-\textcolor{ForestGreen}{\textbf{ongoing}}. \textit{Researcher}. Remote sensing seismic data analysis.

% \href{https://lisa.ulb.ac.be/}{Laboratory of Image Synthesis and Analysis (LISA)}. 
\vspace{5pt}

\href{https://ifz.ru/}{\textbf{Schmidt Institute of Physics of the Earth of the Russian Academy of Sciences}}. Department III: Natural Disasters, Anthropogenic Hazards and Seismicity of the Earth. Laboratory 303: Regional Geophysics and Natural Disasters. Moscow, Russia. 11.2020-\textcolor{ForestGreen}{\textbf{ongoing}}. \textit{Senior Engineer}. Cartographic and and tectonic mapping for geophysical research. State Examinations:  \href{https://drive.google.com/file/d/1D2cxxBIf5pgjPkhkfMwrXsaXrkb3GRPx/view?usp=sharing}{Geotectonics and Geodynamics, English and Philosophy. Certificate}. % Morphostructural features of the deep-sea trenches in the Pacific Ocean, the problem of their origin. 
\vspace{5pt}

\textbf{Analytical Center}. Moscow. 09-10.2020. \textit{Consultant}. Remote sensing data analysis for vegetation indices and land cover types mapping (Landsat TM, Sentinel-2A), ArcGIS, QGIS, SAGA GIS.
\vspace{5pt}

\textbf{\href{http://www.czechglobe.cz/en/}{CzechGlobe}}. \v{C}esk\'{e} Bud\v{e}jovice. 03-08.2014. Division of Ecosystems Analysis. \href{https://www.researchgate.net/project/2013-2015-Czech-Republic-Sumava-National-Park-Environmental-Analysis}{Environmental spatial analysis of the \v{S}umava National Park (QGIS), NATURA-2000}.
\vspace{5pt}

\href{https://www.ntu.edu.tw/english/}{\textbf{National Taiwan University}}. Taipei. 04-05.2013. Department of Geography. \href{https://www.researchgate.net/project/2013-Taiwan-Taipei-Landsat-TM-scenes-and-ENVI-GIS-for-mapping-urban-growth}{ENVI GIS for k-means clustering of Landsat TM/ETM+ for mapping urban sprawl in Taibei since 1990}.
\vspace{5pt}

\textbf{Ege \"{U}niversitesi}. 11-12.2012. Faculty of Geography. Erdas Imagine for GIS analysis of Izmir area, Aegean Sea, Turkey. \href{https://www.researchgate.net/project/2012-Izmir-Turkey-Environmental-Analsysis-of-the-Aegean-Sea-Area}{Land cover change detection by Landsat TM, Google Earth (1990–2010)}
\vspace{5pt}

\href{http://www.ulb.ac.be/}{\textbf{Universit\'e Libre de Bruxelles}}. 09-10.2012. Unit\'e Analyse G\'eospatial. \href{https://www.researchgate.net/project/2012-Belgium-Brussels-OBIA-Object-Based-Image-Analysis-for-Urban-Mapping}{Detecting buildings by OBIA segmentation on Pleiades VHR image using eCognition}.
\vspace{5pt}

\href{http://www.umb.sk/umb/umbbb.nsf}{\textbf{Univerzita Mateja Bela v Banskej Bystrici}}. 07-08.2012. Fakulta prírodných vied. ArcGIS. \href{https://www.researchgate.net/project/2012-Slovakia-Banska-Bystrica-Environmental-Analysis}{Landscape metrics calculation (patch size, density, shape quantification): NATURA 2000}.
\vspace{5pt}

\href{http://www.elte.hu/}{\textbf{E\"otv\"os Lor\'and University}}. 01-04.2012. \href{http://lazarus.elte.hu/}{Dept of Cartography \& GIS}. Mapping land cover changes, Mecsek Hills. \href{https://www.researchgate.net/project/2012-Hungary-Budapest-ILWIS-GIS-for-Monitoring-Central-European-Landscapes}{Land cover change detection by Landsat TM/ETM+: clustering in ILWIS GIS}.
\vspace{5pt}

\href{http://www.arcticcentre.org/InEnglish.iw3}{\textbf{University of Lapland}}. 07-12.2011. Arctic Centre. \href{https://www.researchgate.net/project/2011-Finland-Lapland-ILWIS-GIS-and-Landsat-TM-ETM-for-Mapping-Land-Cover-Changes-in-of-Yamal-Russian-Arctic}{Environmental analysis \& land cover change mapping of Yamal by Landsat TM/ETM+. ILWIS GIS, ArcGIS}
\vspace{5pt}

\href{http://www.uef.fi/tutustu}{\textbf{University of Kuopio (Eastern Finland)}}. 06-09.2009. Department of Environmental Sciences. \href{https://www.researchgate.net/project/2009-Finland-Kuopio-Envionmental-Data-Managmenet}{Data management, organizing and analysis: biochemical data on contaminants in wetlands}
\vspace{5pt}

\href{https://www.awi.de/}{\textbf{Alfred-Wegener-Institut f\"{u}r Polar- und Meeresforschung (AWI})}. 10.2007-03.2008. Dept of Marine Geology and Palaeontology. \href{https://www.researchgate.net/project/2007-2008-Germany-Bremerhaven-AutoTrace-Software-for-Digitizing-Maps}{AutoTrace digitizer for raster2vector conversion}. 
\vspace{5pt}

\href{http://atlas-print.ru/}{\textbf{Atlas-Print}}, retail Department 'Map Trade'. Russia, Moscow. 04-09.2008. \textit{GIS-Specialist}. Updating topographic maps by actual satellite imagery: MapInfo. Map actualization (Moscow city).
\vspace{5pt}

\href{http://www.rniikp.ru/}{\textbf{Russian Research Institute for Space Device Engineering}}. Moscow.12.2006-06.2007. \textit{Research Assistant}. Image processing, geo-referencing, geometric correction, mosaics.
\vspace{5pt}

\href{http://eng.ntsomz.ru/}{\textbf{Research Centre for Earth Operative Monitoring}}. Russia, Moscow. 09-12.2006 \textit{Research Assistant}. Image processing, orthotransformation, overlaying: Erdas Imagine.
\vspace{5pt}

\textbf{Land Agroindustrial Corporation ZNAK}. Russia, Moscow. 08-09.2006. \textit{Specialist in Land Cadastre}. Creating geodatabase for cadastral GIS in GeoMedia.
\vspace{5pt}

\href{https://www.awi.de/}{\textbf{Alfred-Wegener-Institut f\"{u}r Polar- und Meeresforschung (AWI)}}. 01-02.2006 and 04-06.2006. Bathymetry \& Geodesy Research Group. \href{https://www.researchgate.net/project/2006-Antarctica-South-Pacific-Ocean-Bathymetric-Surveying-in-Amundsen-Bellingshausen-Seas-ANT-XXIII-4}{Pre- and post-expedition work: ANT/XXIII-4}.
\vspace{5pt}

\href{http://www.ats.aq/documents/ie/deanu06e.pdf}{\textbf{Antarctic expedition ANT/XXIII-4}}, 02-04.2006. Cartographer onboard the R/V \textit{'Polarstern'}. Sonar data processing: Hydrosweep DS-2. Data conversion, integration (Caris HIPS, ArcGIS). 
\vspace{5pt}

\href{http://www.msu.ru/en/}{\textbf{Lomonosov Moscow State University}}, Faculty of Geography. 09-12.2005. \textit{Engineer}. Thematic mapping, data processing, environmental analysis of the Barents Sea. ArcGIS.
\vspace{5pt}

\textbf{Khimki Municipal Cadastral Centre}. Russia, Khimki. 09.2004-04.2005. \textit{Cartographer}. Editing topographic maps of the Khimki city by AutoCAD Map2005. 
\vspace{5pt}

\textbf{Ltd Svean}. Russia, Moscow. 08.1997-09.1998. Office secretary work: copy, type, scan, print. 

%\pagebreak

\section{Courses}
 
\begin{dinglist}{46}
	\item \href{http://www.padi.com/scuba/}{\textbf{PADI Open Water Diver}}. Crete. \href{http://www.eurodiving.net/}{Eurodiving}. Ligaria, Agia Pelagia, 05.10.2010. \href{https://drive.google.com/file/d/1_w4fjQdQ9nmOj3bsQoDQU2NV-kv12v9L/view?usp=sharing}{Certificate}.\vspace{6pt}
	\item \href{http://www.padi.com/scuba/}{\textbf{PADI SCUBA Water Diver}}. Crete. \href{http://www.eurodiving.net/}{Eurodiving}. Agia Pelagia, 30.09.2010. \href{https://drive.google.com/file/d/19pxzxu13l6IFnRhFRmZA98Ie_MguJSpT/view?usp=sharing}{Certificate}.\vspace{6pt}              
	\item \href{http://www.reson.com/products/seabat/}{SeaBat QTC Multiview and Chesapeake SonarWIZ Echosounder System}. Moscow, 19-26.10.2005. Sonar system \href{http://www.reson.com/products/pds-2000-software/pds2000-multibeam/}{PDS2000}, seafloor mapping. 40 hrs. \href{http://atlantic.ginras.ru/}{Institute of Geology RAS}. \href{https://drive.google.com/file/d/1h9mr-TJyJxDGd5c685OndbsteIjHcPmy/view?usp=sharing}{Certificate}.\vspace{6pt}
	\item \textbf{AutoCAD Map 3D}. Moscow. 2-6/08/2005. Software course 40 hrs. \textit{\href{https://www.pointcad.ru/}{Point Autodesk}}. \href{https://drive.google.com/file/d/13N08X1mpAMH-CE9HrgC3R3ysKAVTf6zN/view?usp=sharing}{Certificate}. \vspace{6pt} 
\end{dinglist}

%\pagebreak

\section{Membership}

\begin{innerlist}
	\item 2013, On Sustainability Knowledge Community
	\item 2013, Heritage \& Archaeological Research Practice
	\item 2012, International Association of Early-Career Geoscientists
	\item 2012, Earth Science Women's Network (ESWN)
	\item 2012, Livestock Emissions \& Abatement Research Network (LEARN)%
	\item 2012, Association for Women Geoscientists (AWG)
	\item 2011, European Geoscience Union (EGU)
	\item 2009, Swiss Snow Ice and Permafrost Society
	\item 2009, Swiss Geomorphological Society (SGS)
\end{innerlist}

\section{Journal Referee}
 
\begin{etaremune} 
%\item \href{}{}. ISSN. \href{}{Publication: }
	\item \href{https://www.frontiersin.org/journals/sustainable-cities#}{Frontiers In Sustainable Cities}. ISSN: 2624-9634.
	\item \href{https://www.frontiersin.org/journals/marine-science}{Frontiers in Marine Science}. ISSN: 2296-7745.
	\item \href{https://www.mdpi.com/journal/remotesensing}{Remote Sensing}. ISSN 2072-4292 
	\item \href{https://journals.plos.org/plosone/}{PLOS One}. ISSN 1932-6203
	\item \href{https://www.tandfonline.com/toc/tgsi20/current}{Geo-spatial Information Science}. ISSN: 1009-5020
	\item \href{https://www.springer.com/journal/10064/}{Bulletin of Engineering Geology and the Environment (Springer Nature)}. ISSN 1435-9529
	\item \href{https://onlinelibrary.wiley.com/journal/17556724}{Acta Geologica Sinica-English Edition}. ISSN 1755-6724.
	\item \href{https://www.tandfonline.com/toc/tgei20/current}{Geocarto International}. ISSN 1010-6049.
	\item \href{https://www.mdpi.com/journal/ijgi}{ISPRS International Journal of Geo-Information}. ISSN 2220-9964. \href{https://drive.google.com/file/d/1ioGWeWVKKqZpwP0HzooXEDj6_Ij4LAER/view?usp=sharing}{Certificate}. \href{https://www.mdpi.com/2220-9964/10/10/667}{Publication: https://doi.org/10.3390/ijgi10100667}
	\item \href{https://www.mdpi.com/journal/ijgi}{ISPRS International Journal of Geo-Information}. ISSN 2220-9964. \href{https://drive.google.com/file/d/1TiqsxF7uSL9UQ21adQilADdoAD4C0Q9v/view?usp=sharing}{Certificate}. \href{https://www.mdpi.com/2220-9964/10/10/647}{Publication: https://www.mdpi.com/2220-9964/10/10/647}
	\item \href{https://www.frontiersin.org/journals/microbiology#}{Frontiers in Microbiology}. ISSN. 1664-302X  \href{https://doi.org/10.3389/fmicb.2021.743920}{Publication: https://doi.org/10.3389/fmicb.2021.743920}
	\item \href{https://www.frontiersin.org/journals/earth-science}{Frontiers in Earth Science}. ISSN 2296-6463. \href{https://doi.org/10.3389/feart.2021.749611}{Publication: https://doi.org/10.3389/feart.2021.749611}
	\item \href{https://www.frontiersin.org/journals/marine-science}{Frontiers in Marine Science}. ISSN: 2296-7745. \href{https://doi.org/10.3389/fmars.2021.730984}{Publication: https://doi.org/10.3389/fmars.2021.730984}
	\item \href{https://www.mdpi.com/journal/land/about}{Land (MDPI)}. ISSN: 2073-445X. \href{https://drive.google.com/file/d/1ezcRdloMhls2jiYzFlFKCJiqBOLxIi4v/view?usp=sharing}{Certificate}. \href{https://doi.org/10.3390/land10080859}{Publication: https://doi.org/10.3390/land10080859}
	\item \href{https://www.tandfonline.com/toc/tgsi20/current}{Geo-spatial Information Science}. ISSN: 1009-5020
	\item \href{https://ceer.com.pl/resources/html/cms/MAINPAGE}{Civil and Environmental Engineering Reports}. ISSN: 2080-5187
	\item \href{https://journals.vgtu.lt/index.php/GAC/about}{Geodesy and Cartography}. ISSN: 2029-6991
	\item \href{https://onlinelibrary.wiley.com/journal/14679671}{Transactions in GIS}. ISSN: 1361-1682
	\item \href{https://www.journals.elsevier.com/progress-in-oceanography}{Progress in Oceanography}. ISSN: 0079-6611. \href{https://drive.google.com/file/d/1su-1xv6EO_PuRypu_Qs3NiCREJlZQk3a/view?usp=sharing}{Certificate}. \href{https://doi.org/10.1016/j.pocean.2020.102477}{Publication: 10.1016/j.pocean.2020.102477}
	\item \href{http://aquatres.scientificwebjournals.com/}{Aquatic Research}. ISSN: 2618-6365. \href{https://doi.org/10.3153/AR19014}{Publication: https://doi.org/10.3153/AR19014}
\end{etaremune} 

\pagebreak
\section{References}

\begin{tabular}{p{2.7cm} p{3cm} p{2.5cm} p{6cm}}  
\toprule
\multicolumn{2}{c}{\textbf{Referee}} \\
\cmidrule(r){1-2}
Name & Title & Country & Organization \\
\midrule
\rowcolor{Lavender}
	\href{https://drive.google.com/file/d/1PUFvQkLrY8ZMwK_pzK_R8R20Dp982U1W/view?usp=sharing}{Michal Klau\v{c}o} & RNDr. PhD. & Slovakia & Matej Bel University, Dept of Geography, Geology \& Landscape Ecology\\
	\href{https://drive.google.com/file/d/1zdeVOFVFPKBZ4024KCJzGKngkC5-IfLM/view?usp=sharing}{B. Gregorová} & PaedDr. PhD. & Slovakia & Matej Bel University, Department of Geography and Geology \\
\rowcolor{Lavender}
	\href{https://drive.google.com/file/d/1Qt-S-5AbajT0SChMXJrifLQg6CRaWuiN/view?usp=sharing}{Thomas Glade} & Prof. Dr. & Austria & University of Vienna, Faculty of Geosciences, Geography and Astronomy \\
	\href{https://drive.google.com/file/d/11-IoaxsmsiHGeDOlpT_oF4ruDVdgSf6I/view?usp=sharing}{Cees van Westen} & Assoc. Prof. Dr. & Netherlands & University of Twente, Faculty of GIS and Earth Observation \\
\rowcolor{Lavender}
	\href{https://drive.google.com/file/d/1_3zIh0jxsPIX9BWuITK4bAbqbXqsx5Um/view?usp=sharing}{A. G. Toxopeus} & Assoc. Prof. Dr. & Netherlands & University of Twente, Faculty of GIS and Earth Observation\\
	\href{https://drive.google.com/file/d/1kSNt05NEZwb8ptl8CAKkQfGuUlt8tUj9/view?usp=sharing}{P. Pilesj\"{o}} & Prof. Dr. Dir. & Sweden & Lund University, GIS Centre, Dept. of Phys. Geogr. and Ecosystem Science \\
\rowcolor{Lavender}
	\href{https://drive.google.com/file/d/1k_S4zYuXzJbQ2lixuqVriK9Pt9m_ck1L/view?usp=sharing}{E. J. Milton} & Emer. Prof. Dr. & UK (England) & University of Southampton,  School of Geography \\
	\href{https://drive.google.com/file/d/1927nFOzndI4IMoP7FH7a-mTL6oUAwciq/view?usp=sharing}{A. Harris} & Sr. Lecturer, Dr. & UK (England) & University of Southampton, School of Geography\\
\rowcolor{Lavender}
	\href{https://drive.google.com/file/d/1oFz8JTVuggPcNsPiRXEVFqg3RuC5S-qQ/view?usp=sharing}{Nikolaus J. Kuhn}  & Prof. Dr. & Switzerland & University of Basel, Dept. of Environmental Sciences \\
	\href{https://drive.google.com/file/d/1mMEDsJv8DA5YfcVBIpF-iSUp6RJINrnh/view?usp=sharing}{H.-W. Schenke} & Hon. Prof. Dr.-Ing.  & Germany & AWI for Polar and Marine Research, Dept. of Geosciences\\
\rowcolor{Lavender}
	\href{https://drive.google.com/file/d/1Sg2aHh6ehWWXJoA3yfZOHRqvFMCQaUMl/view?usp=sharing}{A. M. Berlyant} & Emer. Prof. Dr. & Russia & Lomonosov Moscow State University, Faculty of Geography \\
	\href{https://drive.google.com/file/d/1OJGjFAJMfOn6JbqXfLDDQaW8aSA2fQ_e/view?usp=sharing}{E.A. Prokhorova} & Assoc. Prof. Dr. & Russia & Lomonosov Moscow State University, Faculty of Geography \\
\rowcolor{Lavender}
	\href{https://drive.google.com/file/d/1WxCUw1HLjhMeK00DQX2_ngkvQqOoJZTE/view?usp=sharing}{Y. Wang} & Dr. & China & Ocean University of China, Faculty of Marine Geology (2018) \\
	\href{https://drive.google.com/file/d/1Mi_W7gQ68WxqVHFenCChzE8-0h0D_kAs/view?usp=sharing}{Y. Wang} & Dr. & China & Ocean University of China, Faculty of Marine Geology (2019) \\
\rowcolor{Lavender}
\bottomrule
\end{tabular}

\pagebreak

\section{Publications}
\nocite{*}

\renewcommand{\UrlFont}{\small\tt}

\subsection{Journal Articles}
\newrefcontext[labelprefix=A.]
\printbibliography[type=article, keyword=article, resetnumbers=true, heading=none]

\subsection{Presentations}
\newrefcontext[labelprefix=Pr.]
Presentations at \href{https://www.slideshare.net/PolinaLemenkova/presentations}{SlideShare}.
\printbibliography[type=inproceedings, keyword=presentation, resetnumbers=true, heading=none]

\subsection{Posters}
\newrefcontext[labelprefix=Po.]
Posters at \href{https://www.slideshare.net/PolinaLemenkova/presentations}{SlideShare}.
\printbibliography[type=inproceedings, keyword=poster, resetnumbers=true, heading=none]

%\subsection{Programming Codes}
%\newrefcontext[labelprefix=PC.]
%My codes are publicly available on the \href{https://github.com/paulinelemenkova}{GitHub}.
%\printbibliography[type=misc, keyword=code, resetnumbers=true, heading=none]

\subsection{Conference Proceedings}
\newrefcontext[labelprefix=CP.]
\printbibliography[type=inproceedings, keyword=inproceedings, resetnumbers=true, heading=none]

\subsection{Theses}
\newrefcontext[labelprefix=T.]
\printbibliography[type=thesis, keyword=thesis, heading=none]

%\subsection{Periodicals}
%\newrefcontext[labelprefix=P.]
%\printbibliography[type=periodical, keyword=periodical, resetnumbers=true, heading=none]

\subsection{Technical Reports}
\newrefcontext[labelprefix=TR.]
\printbibliography[type=report, keyword=report, resetnumbers=true, heading=none]

\subsection{Books}
\newrefcontext[labelprefix=B.]
\printbibliography[type=book, keyword = book, resetnumbers=true, heading=none]

\subsection{Book Chapters}
\newrefcontext[labelprefix=BC.]
\printbibliography[type=inbook, keyword = chapter, resetnumbers=true, heading=none]

%\subsection{Datasets}
%\newrefcontext[labelprefix=D.]
%\printbibliography[type=misc, keyword=dataset, resetnumbers=true, heading=none]

%\subsection{Maps}
%\newrefcontext[labelprefix=M.]
%\printbibliography[type=misc, keyword=map, resetnumbers=true, heading=none]

%\subsection{GIS Techniques}
%\newrefcontext[labelprefix=G.]
%\printbibliography[type=misc, keyword=gis, resetnumbers=true, heading=none]

%\subsection{Research Proposals}
%\newrefcontext[labelprefix=RP.]
%\printbibliography[type=misc, keyword=proposal, resetnumbers=true, heading=none]

%\subsection{Student Works}
%\newrefcontext[labelprefix=SW.]
%\printbibliography[type=misc, keyword=student, heading=none]

\end{document}

%%%%%%%%%%%%%%%%%%%%%%%%%% End CV Document of Lemenkova Polina %%%%